
LLM-based code repair systems iteratively refine generated code using feedback from test execution and static analysis. While prior work has studied which feedback types improve repair quality, the effect of feedback ordering under matched content remains unexplored. This work investigates whether presenting static analysis errors before execution failures affects repair outcomes when both conditions receive byte-identical feedback content. Experiments on HumanEval (164 problems) and MBPP (500 problems) with GPT-4o-mini show that static-to-execution ordering achieves 55.57\% pass@1 versus 43.07\% for execution-to-static ordering, a 29.02\% relative improvement (95\% CI: [15.74\%, 44.53\%], p=6.$31\times 10^{-6})$. Mechanism analysis indicates the effect operates through regression prevention: static-first repairs exhibit 38\% fewer inter-iteration regressions (p=0.0198), while early iteration gains are statistically indistinguishable between conditions (p=0.859). These results were obtained using MOCK\_POC execution mode; validation with real API calls is required to confirm exact magnitudes.

